\chapter{Background and Related Work}
\chaptermark{Background and Related Work}
\HRule\\[2cm]
\label{Chapter2}
\lhead{\emph{\textbf{Chapter 2:} Background and Related Work}}

\begin{spacing}{1.6}

\section{Homomorphic Encryption}

\subsection{Historical Overview}

The concept of homomorphic encryption was first described by Rivest, Adleman,
and Dertouzos in 1978 \cite{rivest1978}, who posed the question of whether it
was possible to perform computations on encrypted data. Early schemes such as
RSA and Paillier supported either multiplication or addition homomorphically,
but not both, limiting their utility. The construction of a
\textit{Fully Homomorphic Encryption} (FHE) scheme capable of evaluating
arbitrary circuits was not achieved until Gentry's landmark work in 2009
\cite{gentry2009}, which used ideal lattice-based cryptography. While
theoretically groundbreaking, Gentry's scheme was too computationally expensive
for practical use.

Subsequent work produced a series of increasingly efficient FHE schemes:

\begin{itemize}
  \item \textbf{BGV} (Brakerski, Gentry, Vaikuntanathan, 2012): A Somewhat
  Homomorphic Encryption (SHE) scheme based on Ring Learning With Errors
  (RLWE), supporting integer arithmetic with batch processing via SIMD packing.
  \item \textbf{BFV} (Fan, Vercauteren, 2012): A variant of BGV with improved
  noise management, implemented in Microsoft SEAL.
  \item \textbf{CKKS} (Cheon, Kim, Kim, Song, 2017): The first scheme designed
  for approximate arithmetic on real numbers, enabling efficient computation on
  floating-point data without quantization.
\end{itemize}

\subsection{Mathematical Foundations of CKKS}

The CKKS scheme \cite{ckks} operates over polynomial rings of the form
$\mathcal{R} = \mathbb{Z}[X]/(X^n + 1)$, where $n$ is the polynomial modulus
degree, which must be a power of two. The security of the scheme rests on the
hardness of the Ring Learning With Errors (RLWE) problem.

\textbf{Encoding:} The CKKS encoder maps a vector of $n/2$ complex (or real)
numbers to a polynomial in $\mathcal{R}$ via the inverse of the canonical
embedding. This encoding naturally supports element-wise (SIMD) operations:
adding or multiplying two encoded polynomials corresponds to element-wise
addition or multiplication of the original complex vectors. For a polynomial
modulus degree of $n = 16384$, each ciphertext encodes $8192$ real values
simultaneously.

\textbf{Encryption:} A plaintext polynomial $m \in \mathcal{R}$ is encrypted
by computing:
\begin{equation}
  \mathrm{ct} = (c_0, c_1) = (m + e_0 + pk_0 \cdot u, e_1 + pk_1 \cdot u)
\end{equation}
where $pk = (pk_0, pk_1)$ is the public key, $u$ is a small random polynomial,
and $e_0, e_1$ are small error polynomials drawn from a discrete Gaussian
distribution. The error polynomials provide semantic security under RLWE.

\textbf{Homomorphic Operations:}

\textit{Addition:} For ciphertexts $\mathrm{ct}_1$ and $\mathrm{ct}_2$:
\begin{equation}
  \mathrm{ct}_1 + \mathrm{ct}_2 = (c_0^{(1)} + c_0^{(2)}, c_1^{(1)} + c_1^{(2)})
\end{equation}
Addition is exact (up to encoding precision) and does not consume multiplication
levels.

\textit{Multiplication:} Ciphertext multiplication produces a degree-2
polynomial requiring \textit{relinearization} (using relinearization keys) to
return to the standard form $(c_0, c_1)$. Each multiplication also scales up
the ciphertext, requiring \textit{rescaling} to bring it back to the target
scale:
\begin{equation}
  \mathrm{ct}_1 \otimes \mathrm{ct}_2 \rightarrow
  \mathrm{Relinearize} \rightarrow \mathrm{Rescale}
\end{equation}
Each rescaling operation reduces the coefficient modulus by one prime from the
modulus chain, consuming one multiplication \textit{level}.

\textbf{Level Management:} The coefficient modulus is a product of primes
$q = q_0 \cdot q_1 \cdots q_L$. Each rescaling removes one prime, reducing
the remaining budget by one level. When the budget is exhausted, further
multiplications are not possible without bootstrapping. The number of usable
levels is determined by the length of the modulus chain minus the first and
last special primes.

\textbf{Rotations:} The slot vector can be cyclically shifted using Galois
automorphisms, requiring precomputed Galois keys. A left rotation by $k$
positions maps slot $i$ to slot $i-k$ (mod $n/2$). Rotations do not consume
multiplication levels but do require a key-switching operation.

\textbf{Approximate Arithmetic:} Unlike BFV which is exact for integer
arithmetic, CKKS introduces a bounded approximation error $|\varepsilon|
\leq 2^{-\Delta}$ where $\Delta$ is the scaling factor (typically $2^{40}$
for 40-bit precision). This error is tolerable for ML applications where
exact arithmetic is not required.

\subsection{CKKS Parameters and Their Trade-offs}

\begin{table}[H]
\centering
\caption{CKKS Parameter Definitions and Trade-offs}
\label{tab:ckks_params}
\begin{tabular}{p{3.5cm}p{4cm}p{4cm}}
\toprule
\textbf{Parameter} & \textbf{Effect on Security/Capacity} & \textbf{Effect on Performance} \\
\midrule
\texttt{poly\_modulus\_degree} $n$ &
Larger $n$ $\to$ higher security, more SIMD slots ($n/2$), larger modulus budget &
Larger $n$ $\to$ $O(n \log n)$ NTT operations, higher per-operation cost \\
\midrule
\texttt{coeff\_modulus} chain &
More primes $\to$ more multiplication levels; total bit-length determines security &
Longer chain $\to$ larger ciphertexts, more key-switching cost \\
\midrule
Scale $\Delta = 2^p$ &
Larger $\Delta$ $\to$ more precision, more bits per prime consumed per rescale &
Larger $\Delta$ $\to$ fewer levels available from same modulus budget \\
\midrule
Galois keys &
Required for rotation operations &
Precomputation cost; stored as part of server key material \\
\bottomrule
\end{tabular}
\end{table}

\begin{table}[H]
\centering
\caption{Microsoft SEAL CKKS Configurations Used in This Work}
\label{tab:ckks_config}
\begin{tabular}{lll}
\toprule
\textbf{Parameter} & \textbf{Preprocessing Phase} & \textbf{HE Detection Pipeline} \\
\midrule
Scheme & CKKS & CKKS \\
\texttt{poly\_modulus\_degree} & 8192 & 16384 \\
\texttt{coeff\_modulus} bits & [60, 40, 40, 60] & [60, 40, 40, 40, 40, 40, 60] \\
Scale $\Delta$ & $2^{40}$ & $2^{40}$ \\
SIMD slots & 4096 & 8192 \\
Usable levels & 1 & 5 \\
Security level & 128-bit & 128-bit \\
Decryption error bound & $\leq 10^{-9}$ & $\approx 10^{-9}$ \\
\bottomrule
\end{tabular}
\end{table}

The degree 16384 configuration with 5 levels was selected for the detection
pipeline to accommodate both HE-SAD (2 levels) and HE-PCA (3 levels) with
headroom for future extensions, while maintaining 128-bit security.

\section{Unsupervised Anomaly Detection}

\subsection{Problem Formulation}

Let $\mathcal{X}_\text{train} = \{\mathbf{x}_1, \ldots, \mathbf{x}_N\}$ denote
the training dataset consisting of $N$ observations from the normal class only,
where each observation $\mathbf{x}_i \in \mathbb{R}^d$ is a $d$-dimensional
feature vector. An anomaly detection model learns a function $f: \mathbb{R}^d
\to \mathbb{R}$ that assigns an anomaly score $s = f(\mathbf{x})$ to any test
observation, where higher scores indicate greater deviation from normal
behaviour. A threshold $\tau$ separates normal ($s \leq \tau$) from anomalous
($s > \tau$) observations. The unsupervised setting requires $f$ to be learned
without any labeled anomaly examples.

\subsection{Statistical Anomaly Detection (SAD)}

\subsubsection{Algorithm}

SAD models normal behaviour under a multivariate Gaussian assumption. The
per-feature mean $\mu_j$ and variance $\sigma_j^2$ are estimated from the
normal training data:
\begin{align}
  \mu_j &= \frac{1}{N} \sum_{i=1}^N x_{ij} \\
  \sigma_j^2 &= \frac{1}{N} \sum_{i=1}^N (x_{ij} - \mu_j)^2
\end{align}

The anomaly score for a test observation $\mathbf{x}$ is the diagonal
Mahalanobis distance:
\begin{equation}
  s(\mathbf{x}) = \sum_{j=1}^d \frac{(x_j - \mu_j)^2}{\sigma_j^2}
  \label{eq:sad}
\end{equation}

This score measures the total normalised squared deviation from the normal
mean across all features. Each feature contributes independently, with
features having lower variance (more consistent in normal operation) receiving
proportionally higher weight when they deviate.

\subsubsection{Threshold Computation}

The detection threshold is set according to the three-sigma rule applied to
training scores:
\begin{equation}
  \tau = \bar{s}_\text{train} + 3 \sigma_{s,\text{train}}
\end{equation}
where $\bar{s}_\text{train}$ and $\sigma_{s,\text{train}}$ are the mean and
standard deviation of scores computed on the normal training set. Under the
Gaussian assumption, this threshold captures approximately 99.7\% of the normal
score distribution, yielding a theoretical false positive rate of 0.3\%.

\subsubsection{HE Compatibility}

Every operation in Equation~\eqref{eq:sad} is a polynomial operation:
subtraction ($x_j - \mu_j$), squaring, scalar multiplication (by
$1/\sigma_j^2$), and summation. This makes SAD natively executable under
CKKS encryption without any non-polynomial function approximation.

\subsection{PCA-Based Anomaly Detection}

\subsubsection{Algorithm}

PCA anomaly detection identifies observations that deviate from the normal
operating subspace. Principal Component Analysis is applied to the centred
normal training data to find the directions of maximum variance:
\begin{equation}
  \mathbf{C} = \frac{1}{N-1} \sum_{i=1}^N (\mathbf{x}_i - \boldsymbol{\mu})
  (\mathbf{x}_i - \boldsymbol{\mu})^\top = \mathbf{V} \mathbf{\Lambda}
  \mathbf{V}^\top
\end{equation}
where $\mathbf{V} \in \mathbb{R}^{d \times d}$ contains the eigenvectors
(principal components) and $\mathbf{\Lambda} = \mathrm{diag}(\lambda_1,
\ldots, \lambda_d)$ the corresponding eigenvalues in decreasing order.

The first $k^*$ principal components span the \textit{normal subspace}
$\mathbf{V}_{k^*} \in \mathbb{R}^{k^* \times d}$. The anomaly score is the
reconstruction error when projecting onto this subspace:
\begin{equation}
  s(\mathbf{x}) = \left\| (\mathbf{x} - \boldsymbol{\mu}) -
  \mathbf{V}_{k^*}^\top \mathbf{V}_{k^*} (\mathbf{x} - \boldsymbol{\mu})
  \right\|^2
  \label{eq:pca}
\end{equation}

Normal observations lie near the normal subspace (low reconstruction error)
while anomalous observations deviate from it (high error).

\subsubsection{Optimal Component Selection}

The number of principal components $k^*$ is selected as the minimum $k$ that
explains a prescribed fraction $\theta$ of the total variance:
\begin{equation}
  k^* = \min \left\{ k : \frac{\sum_{i=1}^k \lambda_i}
  {\sum_{i=1}^d \lambda_i} \geq \theta \right\}
  \label{eq:optk}
\end{equation}
A threshold of $\theta = 0.90$ (90\% explained variance) is used in this work,
consistent with standard PCA practice. This selection is performed automatically
by the system and reported in the monitoring dashboard.

\subsubsection{HE Compatibility}

PCA reconstruction error is computable under CKKS because it involves only
inner products (multiply-plain + slot sum), subtraction, and squaring ---
all polynomial operations. No eigendecomposition occurs during inference;
the component matrix $\mathbf{V}_{k^*}$ is fitted in plaintext on the training
data (a model parameter, not individual data) and used as a plaintext constant
during HE scoring.

\subsection{Comparison with Other Methods}

\begin{table}[H]
\centering
\caption{Anomaly Detection Methods: HE Compatibility Analysis}
\label{tab:methods}
\begin{tabular}{p{2.5cm}p{3cm}p{2cm}p{3cm}}
\toprule
\textbf{Method} & \textbf{Anomaly criterion} & \textbf{HE compatible} & \textbf{Reason if not} \\
\midrule
SAD & Per-feature deviation from mean & Yes & Polynomial operations only \\
PCA & Subspace reconstruction error & Yes & Polynomial operations only \\
Isolation Forest & Statistical isolation depth & No & Non-polynomial branching \\
Autoencoder & Neural reconstruction error & Partially & Deep circuits exhaust levels \\
One-class SVM & Distance from hyperplane & No & Kernel evaluation non-polynomial \\
Local Outlier Factor & Local density ratio & No & Non-polynomial neighbourhood \\
\bottomrule
\end{tabular}
\end{table}

SAD and PCA are chosen for this work because they satisfy the polynomial
constraint imposed by CKKS, are fully unsupervised, and have demonstrated
competitive performance on the benchmark datasets used for evaluation.

\section{Related Work}

\subsection{Homomorphic Encryption in Machine Learning}

The application of homomorphic encryption to machine learning inference was
pioneered by CryptoNets \cite{gilad}, which demonstrated convolutional neural
network inference on CKKS-encrypted MNIST data. While achieving 99\% accuracy,
the approach required several minutes per prediction and was limited to simple
polynomial activation functions (squared activations instead of ReLU). The
work established the feasibility of HE-based ML inference but focused entirely
on image classification with public models.

HEAR \cite{hear} (nGraph-HE2) extended HE inference to larger neural networks
using graph-level optimizations to minimize HE depth consumption. Gazelle
\cite{gazelle} combined HE with garbled circuits to handle non-linear
operations efficiently in a two-party protocol. CrypTen \cite{crypten} provides
a secure computing framework combining multiple cryptographic primitives
including HE.

All of these works protect the \textit{query} (the input to the model) while
the model itself is public and held by the server. They do not address the
setting where the data corpus itself is sensitive and must be protected from
the party performing the computation.

\subsection{Privacy-Preserving Anomaly Detection}

Privacy-preserving anomaly detection has been addressed through several
cryptographic approaches:

\textbf{Differential Privacy:} Dwork and Roth \cite{dwork} formalized
differential privacy (DP) as a mechanism for releasing statistical queries
about datasets with bounded privacy loss. DP-based anomaly detection adds
calibrated noise to model parameters or scores to prevent inference of
individual records. While providing strong formal guarantees, the noise addition
degrades detection accuracy, particularly for small anomaly datasets where the
signal-to-noise ratio is already limited.

\textbf{Secure Multi-Party Computation (MPC):} MPC protocols allow multiple
parties to jointly compute a function without revealing their inputs. Applied
to anomaly detection, MPC enables collaborative model training without data
sharing but requires multiple rounds of communication between non-colluding
servers. The communication overhead grows with the complexity of the function
being computed, making it challenging for high-dimensional data.

\textbf{HE-based specific approaches:} Bost et al. \cite{bost2015} proposed
HE-based protocols for classification tasks including naive Bayes and
decision trees but required label exposure for prediction. Carpov et al.
\cite{carpov2019} demonstrated encrypted prediction for healthcare monitoring
but in a supervised setting requiring labeled training data.

\textbf{Gap:} No prior work presents a generalized, fully unsupervised,
end-to-end HE-based anomaly detection system that: (a) protects raw training
and test data from the detection server, (b) requires no labeled anomaly
examples, (c) solves encrypted division exactly, and (d) is deployed as a
practical network application.

\subsection{PIR and HE Systems for Data Retrieval}

The benchmarking study in this thesis was motivated by experience with Private
Information Retrieval (PIR) systems that use HE:

\textbf{Pantheon \cite{pantheon}} enables clients to retrieve values from
key-value stores without revealing which key was queried, using HE and Fermat's
Little Theorem for equality testing.

\textbf{NewtonPIR \cite{newtonPIR}} uses Newton polynomial interpolation within
a single-ciphertext FHE scheme to minimize communication overhead in PIR.

\textbf{FastPIR \cite{fastPIR}} applies RLWE-based encryption with batching
and polynomial packing to achieve sublinear communication cost.

\textbf{Distributed PIR \cite{distributedPIR}} reduces overhead by directing
queries probabilistically to a popular database subset, with fallback to the
full database.

These systems established the performance characteristics of FHE libraries
in a retrieval context, directly informing the library selection benchmarking
study presented in Chapter~3.

\section{Evaluation Datasets}

\subsection{SWaT --- Secure Water Treatment Dataset}

The SWaT dataset \cite{swat} was collected by iTrust, Singapore University of
Technology and Design (SUTD), from a fully operational 6-stage water treatment
testbed over 11 days of continuous operation. The testbed replicates a realistic
water treatment process including ultrafiltration, dechlorination, reverse
osmosis, and ultraviolet dechlorination stages, controlled by a SCADA system
with hierarchical network architecture.

Data was collected under normal operating conditions for 7 days, followed by
4 days of attack experiments designed and executed by cybersecurity researchers.
The attack scenarios cover 41 distinct events including pump manipulation, valve
tampering, sensor spoofing, and coordinated multi-point attacks.

\begin{table}[H]
\centering
\caption{SWaT Dataset Detailed Statistics}
\label{tab:swat_detail}
\begin{tabular}{ll}
\toprule
\textbf{Property} & \textbf{Value} \\
\midrule
Total rows & 1,441,719 \\
Normal rows & 1,387,098 (96.2\%) \\
Attack rows & 54,621 (3.8\%) \\
Total features & 51 \\
Continuous features & 26 (flow rates, levels, pressures) \\
Binary features & 25 (pump/valve states) \\
Sampling rate & 1 second \\
Normal operation period & 7 days \\
Attack period & 4 days \\
Distinct attack scenarios & 41 \\
Training rows used & 5,000 (normal only) \\
Test rows used & 2,000 (1,800 normal + 200 attack) \\
\bottomrule
\end{tabular}
\end{table}

The 51 features include continuous analog sensor readings (flow rates measured
by FIT sensors, water levels by LIT sensors, pressures by PIT sensors) and
binary digital actuator states (pump on/off signals from P-series sensors,
motorized valve open/closed positions from MV-series sensors). This
heterogeneity between continuous and binary features motivates the type-aware
normalization strategy described in Chapter~4.

\subsection{BATADAL --- Battle of Attack Detection Algorithms}

The BATADAL dataset \cite{batadal} was generated from C-Town, a realistic
medium-sized water distribution network benchmark, as part of an international
competition organized to compare anomaly detection algorithms across research
groups. The dataset includes 43 operational parameters covering tank levels,
junction pressures, pump operating states, and flow rates.

\begin{table}[H]
\centering
\caption{BATADAL Dataset Detailed Statistics}
\label{tab:batadal_detail}
\begin{tabular}{ll}
\toprule
\textbf{Property} & \textbf{Value} \\
\midrule
Total features & 43 \\
Continuous features & 28 (tank levels, pressures, flows) \\
Binary features & 15 (pump/valve states) \\
Sampling interval & 1 hour \\
Normal training rows & 8,761 \\
Test rows & 2,089 \\
Attack ratio in test & 19.5\% \\
Label column & \texttt{ATT\_FLAG} (0 = normal, 1 = attack) \\
Attack categories & 7 distinct attack scenarios \\
Attack strategy & SCADA-level concealment \\
\bottomrule
\end{tabular}
\end{table}

A distinctive feature of BATADAL is its attack design philosophy: each attack
scenario was crafted by domain experts to remain statistically plausible at the
individual sensor level, mimicking the appearance of normal operation while
causing physical consequences in the network. This concealment strategy
deliberately confounds per-feature statistical methods (such as SAD, which
analyses each feature independently) while disrupting feature correlations
that PCA captures. This property makes BATADAL a useful stress-test for methods
based on feature correlation structure.

\begin{table}[H]
\centering
\caption{Feature Categorization Comparison}
\label{tab:featcat}
\begin{tabular}{lcc}
\toprule
\textbf{Feature Type} & \textbf{SWaT} & \textbf{BATADAL} \\
\midrule
Continuous & 26 & 28 \\
Binary & 25 & 15 \\
Total & 51 & 43 \\
\bottomrule
\end{tabular}
\end{table}

\end{spacing}
